% =========================================================
% Frequency-Constrained Unit Commitment with
% Multi-Source Virtual Inertia Provision from IBR
%
% Authors:
%   1. Nathaniel Benelisha (corresponding)
%   2. Muhammad Aris Risnandar
%   Supervisor: Ir. Lesnanto Multa Putranto, S.T., M.Eng., Ph.D.,
%               IPM., ASEAN Eng., SMIEEE.
%
% Affiliation: Dept. of Electrical Engineering and Information Technology,
%              Universitas Gadjah Mada, Yogyakarta 55281, Indonesia
%
% Version : v2.0 — September 2026
%           (With actual simulation results from 7 CPLEX scenarios)
% Target  : IEEE Access / Energies (MDPI)
% =========================================================

\documentclass[journal]{IEEEtran}

% --- Packages ---
\usepackage{amsmath,amssymb,amsfonts}
\usepackage{graphicx}
\usepackage{booktabs}
\usepackage{cite}
\usepackage{xcolor}
\usepackage{multirow}
\usepackage{tabularx}
\usepackage{array}
\usepackage{url}

% --- Prevent formula overflow ---
\setlength{\arraycolsep}{2pt}
\allowdisplaybreaks

\begin{document}

% =========================================================
\title{Frequency-Constrained Unit Commitment with\\
Multi-Source Virtual Inertia Provision\\
from Inverter-Based Resources}

\author{Nathaniel~Benelisha,~Muhammad~Aris~Risnandar,~and~%
Ir.~Lesnanto~Multa~Putranto,~S.T.,~M.Eng.,~Ph.D.,~IPM.,~ASEAN~Eng.,~\IEEEmembership{SMIEEE}%
\thanks{N.~Benelisha and M.~A.~Risnandar are graduate students at the
Department of Electrical Engineering and Information Technology,
Universitas Gadjah Mada (UGM), Yogyakarta 55281, Indonesia.
E-mail: nathanielbenelisha@mail.ugm.ac.id}%
\thanks{M.~A.~Risnandar (NIM: 25/571656/STK/01375) is a Doctoral student
at the Program Studi Doktor Teknik Elektro, UGM.}%
\thanks{Ir.~L.~M.~Putranto is a faculty member (Lektor Kepala / Associate Professor)
at the same department. He is the supervising lecturer.
E-mail: lesnanto@ugm.ac.id}}

\maketitle

% =========================================================
\begin{abstract}
The rapid proliferation of Inverter-Based Resources (IBR)---including
solar photovoltaic (PV), wind turbines (WT), and Battery Energy
Storage Systems (BESS)---reduces natural rotational inertia in modern
power systems, causing higher Rate of Change of Frequency (RoCoF)
and deeper frequency nadirs following N-1 generation contingencies.
This paper proposes an enhanced \textit{Frequency-Constrained Unit
Commitment} (FCUC) model that simultaneously co-optimizes conventional
generator scheduling and virtual inertia provision from multiple IBR
sources within a unified Mixed-Integer Linear Programming (MILP)
framework implemented in Python 3.9/Pyomo and solved with IBM
CPLEX~22.1.1.
Seven simulation scenarios on a modified IEEE 10-unit test system
with RoCoF limit of 0.55~Hz/s are evaluated.
The classical UC (Scenario~1) yields a total cost of
\$4{,}029{,}223/day but violates the RoCoF limit in all 24 scheduling
periods (worst-case 0.903~Hz/s).
Enforcing a minimum inertia constraint via synchronous generators
only (Scenario~4) restores frequency security at a cost penalty
of $+$3.65\% (\$4{,}176{,}358/day).
The proposed multi-source virtual inertia model (Scenario~7,
adding BESS-VI and WT-VI) reduces this cost to \$3{,}939{,}091/day,
a saving of 5.69\% versus Scenario~4 and 2.24\% below even the
unconstrained classical UC, while maintaining full RoCoF compliance
with an average system inertia of 43{,}256~MWs.
These results demonstrate that co-optimizing renewable integration
and virtual inertia provision can simultaneously reduce operating
costs and ensure frequency security---a key insight for Indonesia's
RUPTL 2025--2034 target of 61\% renewable generation.
\end{abstract}

\begin{IEEEkeywords}
Frequency-Constrained Unit Commitment (FCUC), Virtual Inertia,
Inverter-Based Resources (IBR), Battery Energy Storage System (BESS),
Wind Turbine, RoCoF, Frequency Nadir, MILP, Pyomo, CPLEX,
Power System Stability, Indonesia, RUPTL 2025--2034.
\end{IEEEkeywords}

% =========================================================
\section{Introduction}
\label{sec:intro}

\subsection{Motivation and Background}

Indonesia's State Electricity Company (PLN) National Electricity
Supply Business Plan (RUPTL) 2025--2034 targets 69.5~GW of new
generation capacity, with 61\% (${\approx}42.4$~GW) from renewable
energy sources~\cite{PLN2025}. This mirrors a global trend in which
inverter-interfaced generation increasingly displaces conventional
synchronous machines (SG), introducing the \emph{inertia deficit
problem}~\cite{Tielens2016}.

Unlike SGs, which store kinetic energy in rotating masses and
release it instantaneously to arrest frequency deviations, IBR
are decoupled from the grid via power-electronic converters and
contribute no natural inertia. As IBR penetration grows, total
system inertia decreases, causing:

\begin{enumerate}
  \item \textbf{Higher RoCoF}: Inversely proportional to total
    system inertia $H_{\text{sys}}$; can trigger protective relay
    operation and cascading failures~\cite{Tielens2016};
  \item \textbf{Deeper frequency nadir}: The minimum frequency
    reached may fall below UFLS thresholds (49.0~Hz in Indonesia);
  \item \textbf{Degraded quasi-steady-state (QSS) frequency}: If
    insufficient governor capacity is committed, the settled
    post-disturbance frequency may be inadequate.
\end{enumerate}

These risks have been demonstrated in major incidents: the 2016
South Australia blackout and the 2019 Great Britain blackout both
resulted from low-inertia conditions~\cite{AEMO2017}.
Indonesia's isolated grid systems (Sulawesi, Kalimantan, Maluku,
Papua) with rapidly growing renewable portfolios are increasingly
exposed to similar risks~\cite{E02Sulawesi,E03JAMALI}.

\subsection{Literature Review}

\subsubsection{Unit Commitment --- Definition and History}

Unit Commitment (UC) is a fundamental optimization problem in power
system operations that determines the optimal on/off schedule
(\textit{commitment}) of generating units over a planning horizon
--- typically 24 to 168 hours --- to minimize total operating costs
while satisfying load demand and operational
constraints~\cite{Padhy2004}. Formally, UC is a large-scale
combinatorial optimization problem involving binary decision variables
(unit status $u_{i,t} \in \{0,1\}$), continuous variables (power
output $P_{i,t}$), and constraints including power balance, generation
capacity limits, ramp rates, minimum up/down times, and spinning
reserve requirements~\cite{Conejo2010}.

The UC problem was first recognized through the work of
Baldwin \textit{et al.} (1959) on economic shutdown of generating
units, marking the transition from manual scheduling toward
systematic optimization~\cite{Padhy2004}. Since then, UC has become
one of the most extensively studied optimization problems in power
engineering, with hundreds of publications comprehensively
categorized in the bibliographical survey by Padhy~\cite{Padhy2004}.

\subsubsection{Evolution of UC Solution Methods}

Solution techniques for UC have evolved significantly over the
past six decades~\cite{Padhy2004}.

\textbf{Priority List (PL)} is the earliest and simplest approach:
units are ranked by fuel cost efficiency and committed sequentially
until demand and reserve requirements are met. While computationally
fast, PL often yields suboptimal solutions because it ignores
inter-temporal coupling constraints.

\textbf{Dynamic Programming (DP)} provided the first rigorous
mathematical framework, decomposing the combinatorial UC into
smaller staged sub-problems. However, DP suffers from the
\textit{curse of dimensionality} --- computational complexity
grows exponentially with the number of generating units, making
it impractical for large-scale systems.

\textbf{Lagrangian Relaxation (LR)} became the historically
dominant method for large-scale UC. By relaxing binding constraints
(e.g., power balance) into the objective function via Lagrange
multipliers, LR decomposes the problem into per-unit sub-problems
that are individually tractable, providing a lower bound on the
optimal cost.

\textbf{Mixed-Integer Linear Programming (MILP)} is now the
dominant approach in modern power system operations, enabled by
advances in commercial solvers such as IBM CPLEX and
Gurobi~\cite{Carrion2006}. Carri\'on and Arroyo~\cite{Carrion2006}
proposed a computationally efficient MILP formulation for thermal
UC requiring fewer binary variables and constraints than earlier
methods; this formulation became the standard benchmark.
MILP models unit on/off decisions as binary variables and
operational constraints as linear inequalities, guaranteeing
global optimality within a specified tolerance gap.
Conejo \textit{et al.}~\cite{Conejo2010} further extended this
framework for decision-making under uncertainty in electricity
markets, including stochastic UC formulations.

\subsubsection{Low-Inertia Power System Challenges}

Rotational inertia is a fundamental physical property of
synchronous generators (SG): kinetic energy stored in the spinning
rotor mass provides an instantaneous frequency arresting response
following power imbalances~\cite{Kundur1994}. As reviewed by
Tielens and Van~Hertem~\cite{Tielens2016}, when SGs are displaced
by IBR --- solar PV, wind turbines, and BESS --- system inertia
declines proportionally because IBR are decoupled from the grid
via power-electronic converters.

The consequences of reduced inertia are characterized by the
System Frequency Response (SFR) model~\cite{Anderson1990}, a
second-order transfer function approximating the closed-loop
frequency dynamics. Anderson and Mirheydar~\cite{Anderson1990}
derived three key frequency metrics from this model:
(i)~RoCoF $= f_0 \Delta P_{\text{loss}} / (2 H_{\text{sys}})$,
inversely proportional to system inertia;
(ii)~frequency nadir, the maximum deviation; and
(iii)~quasi-steady-state (QSS) deviation after primary response
completes.

The severity of low-inertia risks has been demonstrated in
major incidents: the 2016 South Australia blackout (loss of
850~MW; AEMO identified low inertia as a key contributing
factor~\cite{AEMO2017}) and the 2019 Great Britain blackout
(1{,}131~MW loss causing frequency to fall to 48.8~Hz, triggering
UFLS disconnecting ${>}$1~million customers).

\subsubsection{Frequency-Constrained Unit Commitment}

The FCUC concept addresses the fundamental limitation of classical
UC: cost-optimal schedules produced without frequency dynamics
consideration can be \textit{frequency-insecure} when system inertia
is low. FCUC incorporates frequency security constraints (RoCoF,
nadir, QSS) directly into the optimization formulation.

The development of FCUC has progressed chronologically as follows:

\textbf{Ahmadi and Ghasemi (2014)}~\cite{Ahmadi2014} formulated
security-constrained UC with linearized system frequency limit
constraints, demonstrating that linear approximations of the nadir
constraint enable solution via standard MILP solvers.

\textbf{Teng, Trovato, and Strbac (2016)}~\cite{Teng2016} proposed
stochastic scheduling with inertia-dependent fast frequency response
(FFR) requirements. Their model showed that FFR requirements increase
non-linearly as system inertia decreases, establishing the
relationship between scheduling decisions and frequency security.

\textbf{Wen, Li, Huang, and Liu (2016)}~\cite{Wen2016} proposed the
first MILP-based FCUC explicitly embedding RoCoF and nadir
constraints derived from the SFR model~\cite{Anderson1990}, with
BESS as a fast-acting frequency responder. This seminal paper
established the framework for integrating BESS into FCUC.

\textbf{Xu, Jang, and Overbye (2018)}~\cite{Xu2018} extended this
concept with commitment of fast-responding storage devices to mimic
inertia, enhancing primary frequency response without requiring
additional synchronous generator commitment.

\textbf{Chu, Markovic, Hug, and Teng (2020)}~\cite{Chu2020} proposed
a stochastic UC model incorporating synthetic inertia provision
from wind turbines via a two-stage control framework. Their key
innovation was treating wind turbine synthetic inertia as a
\textit{dispatchable ancillary service} and eliminating the
\textit{secondary frequency dip} problem through careful rotor
speed recovery management.

\textbf{Asiedu \textit{et al.} (2025)}~\cite{Asiedu2025} provided
a comprehensive review of dynamic constraint integration into
steady-state optimization models for IBR-dominated systems,
identifying key research gaps and open problems.

\subsubsection{Virtual Inertia: Concepts and Technologies}

Virtual inertia (VI) is a control strategy implemented in
power-electronic converters to emulate the inertial response and
damping characteristics of a synchronous generator. By integrating
energy storage (batteries, DC-link capacitors) with specialized
control loops, grid-forming inverters can autonomously inject or
absorb power during frequency deviations, slowing RoCoF and
improving nadir.

\textbf{BESS Virtual Inertia.} BESS provides VI through
grid-forming converter control: the converter responds to frequency
deviations by injecting/absorbing power from the battery. BESS
offers sub-millisecond response time and bidirectional operation
capability. Optimal placement of ESS-based VI for frequency
stability has been studied by Fang \textit{et al.}~\cite{VI06OptimalPlacement}.

\textbf{Wind Turbine Virtual Inertia.} Modern wind turbines (DFIG
and PMSG types) store kinetic energy in the rotating rotor mass.
Synthetic inertia is obtained by transiently extracting rotor
kinetic energy when frequency drops, then gradually recovering
rotor speed. The primary challenge is the \textit{secondary
frequency dip} that occurs during rotor speed recovery;
Chu \textit{et al.}~\cite{Chu2020} addressed this through their
two-stage control framework.

\subsubsection{Linearization of Frequency Nadir Constraint}

The frequency nadir constraint is inherently nonlinear --- it is a
function of $H_{\text{sys}}$, $R_{\text{sys}}$,
$\Delta P_{\text{loss}}$, and damping parameters. To maintain MILP
tractability, linearization is essential.

Shi \textit{et al.}~\cite{Shi2025} proposed a two-stage
linearization method: Stage~1 maps commitment decisions linearly
to system frequency parameters ($H_{\text{sys}}, R_{\text{sys}}$);
Stage~2 approximates the nadir as a linear function of these
parameters via multivariate regression. This approach enables
tractable MILP formulation with high accuracy.
Liu \textit{et al.}~\cite{Liu2024} further addressed the nadir
constraint for high-penetration island power systems where
system inertia fluctuates significantly.

\subsubsection{Indonesian Power System Context}

Indonesia's RUPTL 2025--2034~\cite{PLN2025} targets 61\%
renewable generation, requiring massive IBR integration that will
significantly reduce system inertia --- particularly in isolated
grids (Sulawesi, Kalimantan, Maluku, Papua).

UC with primary frequency regulation in the Southern Sulawesi
isolated grid~\cite{E02Sulawesi} demonstrated a 7.5\% cost
increase when PFR constraints are enforced, but enabled BESS to
extend VRE hosting capacity to 36--41\% of peak load.
A complementary study on the Jawa-Madura-Bali (JAMALI)
system~\cite{E03JAMALI} incorporated a 2-minute FFR constraint,
requiring 675~MW FFR with only a 0.3\% cost increase.
Nguyen-Hong and Nakanishi~\cite{D02ForecastError} examined FCUC
considering BESS and renewable forecast error, highlighting the
importance of uncertainty in UC scheduling.

\subsubsection{Research Gaps}

Despite the extensive body of literature reviewed above, three
critical research gaps remain:
\begin{itemize}
  \item \textbf{(G1)} Existing models consider BESS \emph{or}
    wind turbines as virtual inertia providers, but not
    \emph{simultaneous multi-source provision} from BESS and
    wind turbines within a unified MILP framework;
  \item \textbf{(G2)} Quantitative cost--security trade-off
    analysis across progressively more constrained FCUC
    configurations --- from classical UC to full multi-source
    VI FCUC --- is not available in the literature;
  \item \textbf{(G3)} Comprehensive case studies combining
    multi-source IBR virtual inertia aligned to national energy
    targets in a rapidly electrifying developing economy
    (such as Indonesia's RUPTL 2025--2034) are lacking.
\end{itemize}

\subsection{Contributions}

This paper addresses the identified gaps through:
\begin{enumerate}
  \item \textbf{Multi-Source VI FCUC Formulation:} A unified MILP
    FCUC model co-optimizing virtual inertia from BESS and wind
    turbines simultaneously with conventional generator scheduling,
    implemented in Python~3.9/Pyomo and solved with IBM CPLEX~22.1.1;

  \item \textbf{Seven-Scenario Systematic Analysis:} Evaluation across
    seven constraint configurations (Table~\ref{tab:scenarios})
    quantifies the cost--security trade-off, yielding actual
    numerical results with all seven scenarios fully executed;

  \item \textbf{Practical Insights for Indonesia:} Case study
    parameters aligned with Indonesian power systems and RUPTL
    2025--2034 targets, demonstrating that full frequency security
    can be achieved below the baseline classical UC cost.
\end{enumerate}

% =========================================================
\section{System Modeling}
\label{sec:model}

\subsection{Conventional Generator Model}

The fuel cost of unit~$i$ at period~$t$ is approximated by piecewise
linear (PWL) cost with $K$ segments to maintain MILP tractability:
\begin{equation}
  C_i^{\text{gen}}(P_{i,t}) \approx
    c_i^{\text{fixed}}\,u_{i,t}
    + \sum_{k=1}^{K} \lambda_{i,k}\,\Delta P_{i,k,t}
  \label{eq:pwl}
\end{equation}
where $\lambda_{i,k}$~[USD/MWh] is the slope of segment~$k$,
$\Delta P_{i,k,t}$~[MW] is power allocated to that segment,
and $u_{i,t}\in\{0,1\}$ is the binary commitment status.
Each committed SG contributes rotational inertia constant
$H_{\text{SG},i}$~[MWs/MVA] (typically 2--9~MWs/MVA~\cite{Kundur1994}).

\subsection{Battery Energy Storage System (BESS) Model}

The State of Charge (SoC) of BESS unit~$b$ evolves as:
\begin{equation}
  E_{b,t} = E_{b,t-1}
    + \eta_c^b\,P_{c,b,t}\,\Delta t
    - \frac{P_{d,b,t}\,\Delta t}{\eta_d^b}
  \label{eq:soc}
\end{equation}
where $E_{b,t}$~[MWh] is stored energy, $P_{c,b,t}$ and
$P_{d,b,t}$~[MW] are charging and discharging power, and
$\eta_c^b$, $\eta_d^b$ are charge/discharge efficiencies.

BESS provides virtual inertia via grid-forming converter control.
The equivalent virtual inertia constant contributed by BESS unit~$b$:
\begin{equation}
  H_{\text{VI,BESS},b,t}
    = K_{\text{VI}}^{\text{BESS}}\cdot P_{\text{VI},b,t}
  \label{eq:bess_vi}
\end{equation}
where $K_{\text{VI}}^{\text{BESS}}$~[MWs/MW] is the virtual inertia
gain coefficient and $P_{\text{VI},b,t}$~[MW] is the power headroom
reserved for VI provision~\cite{VI06OptimalPlacement}.

\subsection{Wind Turbine Model}

Synthetic inertia from wind turbines is provided by transiently
extracting kinetic energy from the rotating rotor mass following
the two-stage control framework of~\cite{Chu2020}:
\begin{equation}
  H_{\text{VI,WT},w,t}
    = K_{\text{VI}}^{\text{WT}}\cdot P_{\text{VI,WT},w,t}
  \label{eq:wt_vi}
\end{equation}
where $K_{\text{VI}}^{\text{WT}}$~[MWs/MW] is the WT virtual
inertia gain.
VI is constrained to the available curtailment headroom:
\begin{equation}
  P_{\text{VI,WT},w,t} \leq P_{\text{Curt},w,t}
  \label{eq:wt_vi_curtailment}
\end{equation}

\subsection{Aggregate System Inertia}

Total effective system inertia at scheduling period~$t$:
\begin{equation}
  H_{\text{sys}}(t)
    = \!\sum_{i \in \mathcal{G}_{\text{th}}}\!
        H_{\text{SG},i}\,P_i^{\max}\,u_{i,t}
      + \!\sum_{b}\!
        K_{\text{VI}}^{\text{BESS}}\,P_{\text{VI},b,t}
      + \!\sum_{w}\!
        K_{\text{VI}}^{\text{WT}}\,P_{\text{VI,WT},w,t}
  \label{eq:hsys}
\end{equation}
This expression is \emph{linear} in the decision variables
$u_{i,t}$, $P_{\text{VI},b,t}$, and $P_{\text{VI,WT},w,t}$,
preserving MILP tractability.

\subsection{System Frequency Response (SFR) Model}

The aggregate frequency response following generation
loss~$\Delta P_{\text{loss}}$ is described by the swing
equation~\cite{Anderson1990}:
\begin{equation}
  2\,H_{\text{sys}}(t)\,\frac{d\Delta f}{dt}
    = -\Delta P_{\text{loss}}
      + \Delta P_{\text{gov}}(t)
      + \Delta P_{\text{VI}}(t)
      - D\,\Delta f(t)
  \label{eq:swing}
\end{equation}
where $\Delta f = f - f_0$ is the frequency deviation from nominal
($f_0 = 50$~Hz), $\Delta P_{\text{gov}}(t)$ is the aggregate
governor response, $\Delta P_{\text{VI}}(t)$ is total virtual
inertia power injection, and $D$ is the load damping coefficient.

From the SFR model, three key frequency metrics are derived.

\textbf{RoCoF} (at $t=0^+$):
\begin{equation}
  \text{RoCoF}
    = \frac{f_0\cdot\Delta P_{\text{loss}}}{2\,H_{\text{sys}}(t)}
  \quad [\text{Hz/s}]
  \label{eq:rocof}
\end{equation}

\textbf{QSS frequency deviation}:
\begin{equation}
  \Delta f_{\text{QSS}}
    = \frac{f_0\cdot\Delta P_{\text{loss}}}
           {D\,S_{\text{sys}} + R_{\text{sys}}(t)}
  \label{eq:qss}
\end{equation}
where $R_{\text{sys}}(t) = \sum_{i}\,R_i\,u_{i,t}$~[MW/Hz]
is the total governor droop response.

% =========================================================
\section{FCUC Problem Formulation}
\label{sec:formulation}

\subsection{Objective Function}

\begin{equation}
  \min \sum_{t=1}^{T} \biggl[
    \sum_{i \in \mathcal{G}} \bigl(
      C_i^{\text{gen}}(P_{i,t})
      + C_i^{\text{SU}}\,v_{i,t}
      + C_i^{\text{SD}}\,w_{i,t}
    \bigr)
    + c_{\text{deg}}\!\sum_{b}
      (P_{c,b,t}+P_{d,b,t})\,\Delta t
  \biggr]
  \label{eq:obj}
\end{equation}
where $v_{i,t}, w_{i,t}\in\{0,1\}$ are start-up/shut-down binary
variables, $C_i^{\text{SU}},C_i^{\text{SD}}$~[USD] are corresponding
costs, and $c_{\text{deg}}$~[USD/MWh] is BESS degradation cost.

\subsection{Standard UC Constraints}

\textbf{C1 --- Power balance}:
\begin{equation}
  \sum_{i \in \mathcal{G}} P_{i,t}
  + \sum_{b}(P_{d,b,t} - P_{c,b,t})
  + \sum_{w} P_{\text{WT},w,t}
  = D_t \quad \forall\,t
  \label{eq:balance}
\end{equation}

\textbf{C2 --- Generation capacity limits}:
\begin{equation}
  u_{i,t}\,P_i^{\min} \leq P_{i,t} \leq u_{i,t}\,P_i^{\max}
  \quad \forall\,i,\,t
  \label{eq:genlim}
\end{equation}

\textbf{C3 --- Ramp rate limits}:
\begin{align}
  P_{i,t} - P_{i,t-1}
    &\leq RU_i\,u_{i,t-1} + P_i^{\min}\,v_{i,t}
  \label{eq:rampup} \\
  P_{i,t-1} - P_{i,t}
    &\leq RD_i\,u_{i,t} + P_i^{\min}\,w_{i,t}
  \label{eq:rampdown}
\end{align}

\textbf{C4 --- Minimum up/down time}:
\begin{align}
  \sum_{\tau=t}^{t+UT_i-1} u_{i,\tau}
    &\geq UT_i\cdot v_{i,t}
  \label{eq:mut} \\
  \sum_{\tau=t}^{t+DT_i-1} (1-u_{i,\tau})
    &\geq DT_i\cdot w_{i,t}
  \label{eq:mdt}
\end{align}

\textbf{C5 --- Start-up/Shut-down logic}:
\begin{equation}
  v_{i,t} - w_{i,t} = u_{i,t} - u_{i,t-1}
  \quad \forall\,i,\,t
  \label{eq:susd}
\end{equation}

\textbf{C6 --- N-1 Spinning reserve}:
\begin{align}
  \Delta P_{\text{loss}}(t)
    &\geq P_{i,t}
  \quad \forall\,i \in \mathcal{G}_{\text{th}},\,t
  \label{eq:largestloss} \\
  \sum_{i \in \mathcal{G}} R_{i,t}
    &\geq \Delta P_{\text{loss}}(t)
  \quad \forall\,t
  \label{eq:reserve}
\end{align}

\subsection{BESS Operational Constraints}

\textbf{C7 --- SoC dynamics and bounds}:
\begin{align}
  E_b^{\min}
    &\leq E_{b,t} \leq E_b^{\max}
    \quad \forall\,b,\,t
  \label{eq:socbound} \\
  E_{b,T} &= E_{b,0}
  \quad \forall\,b \quad \text{(daily energy neutrality)}
  \label{eq:energyneutral}
\end{align}

\textbf{C8 --- BESS power limits with VI headroom}:
\begin{align}
  0 \leq P_{c,b,t}
    &\leq P_b^{\max}\,\delta_{c,b,t}
  \label{eq:chg} \\
  0 \leq P_{d,b,t}
    &\leq P_b^{\max}\,\delta_{d,b,t}
  \label{eq:dis} \\
  \delta_{c,b,t} + \delta_{d,b,t}
    &\leq 1
  \label{eq:nosimcd}
\end{align}

\textbf{C9 --- VI power headroom}:
\begin{equation}
  0 \leq P_{\text{VI},b,t}
    \leq P_b^{\max}\,\delta_{d,b,t} - P_{d,b,t}
  \label{eq:vi_headroom}
\end{equation}

\subsection{Wind Turbine Constraints}

\textbf{C10 --- WT output with curtailment}:
\begin{align}
  P_{\text{WT},w,t}
    + P_{\text{Curt},w,t}
    &= P_{\text{WT},w,t}^{\text{avail}}
  \label{eq:wt_avail} \\
  P_{\text{WT},w,t},\,P_{\text{Curt},w,t}
    &\geq 0
  \label{eq:wt_nn}
\end{align}

\subsection{Frequency Security Constraints}

\textbf{C11 --- RoCoF constraint}
(rearranging~\eqref{eq:rocof} with limit
$\text{RoCoF}^{\max}$~[Hz/s]):
\begin{equation}
  \boxed{
  2\,\text{RoCoF}^{\max}\cdot H_{\text{sys}}(t)
  \geq f_0\cdot\Delta P_{\text{loss}}(t)
  }
  \quad \forall\,t
  \label{eq:rocof_con}
\end{equation}
This is \emph{linear} in $H_{\text{sys}}(t)$ via~\eqref{eq:hsys}.

\textbf{C12 --- Minimum system inertia}:
\begin{equation}
  H_{\text{sys}}(t) \geq H_{\text{sys}}^{\min}
  \quad \forall\,t
  \label{eq:hmin}
\end{equation}

\textbf{C13 --- QSS constraint}
(rearranging~\eqref{eq:qss}):
\begin{equation}
  R_{\text{sys}}(t)
    \geq \frac{f_0\cdot\Delta P_{\text{loss}}(t)}
              {\Delta f_{\text{QSS}}^{\max}}
         - D\,S_{\text{sys}}
  \quad \forall\,t
  \label{eq:qss_con}
\end{equation}
where $\Delta f_{\text{QSS}}^{\max} = f_0 - f_{\text{QSS}}^{\min}$~[Hz].

The complete FCUC model is a \textbf{MILP}:
\begin{itemize}
  \item Decision variables:
    $\{u_{i,t},\,v_{i,t},\,w_{i,t},\,P_{i,t},\,E_{b,t},\,P_{c,b,t},\,
     P_{d,b,t},\,\delta_{c,b,t},\,\delta_{d,b,t},\,P_{\text{WT},w,t},\,
     P_{\text{Curt},w,t},\,P_{\text{VI},b,t},\,P_{\text{VI,WT},w,t}\}$
  \item Objective: minimize~\eqref{eq:obj}
  \item Constraints: C1--C13 (\eqref{eq:balance}--\eqref{eq:qss_con})
\end{itemize}

% =========================================================
\section{Case Study Setup}
\label{sec:caseStudy}

\subsection{Test System}

The model is validated on a \textbf{modified IEEE 10-unit test
system}~\cite{Kazarlis1996} with added IBR (BESS and wind turbine).
Generator parameters are in Table~\ref{tab:gen} and IBR/frequency
parameters in Table~\ref{tab:ibr}.

\begin{table}[!t]
\renewcommand{\arraystretch}{1.2}
\caption{Conventional Generator Parameters}
\label{tab:gen}
\centering\footnotesize
\begin{tabular}{crrrrrrr}
\toprule
Unit & $P^{\min}$ & $P^{\max}$ & $b$ & $c^{\text{fix}}$
     & $H$ & $R_i$ & $UT$/$DT$\\
     & [MW] & [MW] & [\$/MWh] & [\$/h]
     & [MWs/ & [MW/ & [h]\\
     & & & & & MVA] & Hz] &\\
\midrule
G1  & 150 & 455 & 16.19 & 1000 & 5.0 & 200 & 8/8 \\
G2  & 150 & 455 & 17.26 &  970 & 5.0 & 200 & 8/8 \\
G3  &  20 & 130 & 16.60 &  700 & 4.0 &  60 & 5/5 \\
G4  &  20 & 130 & 16.50 &  680 & 4.0 &  60 & 5/5 \\
G5  &  25 & 162 & 19.70 &  450 & 3.5 &  75 & 6/6 \\
G6  &  20 &  80 & 22.26 &  370 & 3.5 &  40 & 3/3 \\
G7  &  25 &  85 & 27.74 &  480 & 3.0 &  40 & 3/3 \\
G8  &  10 &  55 & 25.92 &  660 & 2.5 &  25 & 1/1 \\
G9  &  10 &  55 & 27.27 &  665 & 2.5 &  25 & 1/1 \\
G10 &  10 &  55 & 27.79 &  670 & 2.5 &  25 & 1/1 \\
\bottomrule
\end{tabular}
\end{table}

\begin{table}[!t]
\renewcommand{\arraystretch}{1.2}
\caption{IBR, BESS, and Frequency Security Parameters}
\label{tab:ibr}
\centering\footnotesize
\begin{tabular}{lrr}
\toprule
Parameter & Value & Unit \\
\midrule
BESS rated power ($P_{b}^{\max}$)              & 200   & MW      \\
BESS rated energy ($E_{b}^{\max}$)             & 400   & MWh     \\
BESS SoC limits ($E_{b}^{\min}/E_{b}^{\max}$) & 10/90 & \%      \\
BESS charge/discharge efficiency ($\eta_c,\eta_d$) & 0.95 & ---  \\
BESS VI gain ($K_{\text{VI}}^{\text{BESS}}$)  & 10    & MWs/MW  \\
Wind turbine capacity (WT1)                    & 400   & MW      \\
WT VI gain ($K_{\text{VI}}^{\text{WT}}$)       & 4     & MWs/MW  \\
Nominal frequency ($f_0$)                      & 50    & Hz      \\
RoCoF limit ($\text{RoCoF}^{\max}$)           & 0.55  & Hz/s    \\
Minimum nadir ($f_{\text{nadir}}^{\min}$)      & 49.0  & Hz      \\
Minimum QSS frequency ($f_{\text{QSS}}^{\min}$) & 49.5 & Hz    \\
MIP gap tolerance                              & 0.1   & \%      \\
Solver time limit                              & 300   & s       \\
\bottomrule
\end{tabular}
\end{table}

\subsection{Simulation Scenarios}

Seven scenarios are defined from the modular \texttt{CONFIG}
dictionary in the Python implementation.
Table~\ref{tab:scenarios} shows the constraint activation matrix
(matching the presentation on 23 June 2026).

\begin{table}[!t]
\renewcommand{\arraystretch}{1.2}
\caption{Simulation Scenario Constraint Activation Matrix}
\label{tab:scenarios}
\centering\footnotesize
\begin{tabular}{lccccccc}
\toprule
\textbf{Constraint} & \textbf{S1} & \textbf{S2} & \textbf{S3}
  & \textbf{S4} & \textbf{S5} & \textbf{S6} & \textbf{S7} \\
\midrule
Power balance, C1--C5    & \checkmark & \checkmark & \checkmark
  & \checkmark & \checkmark & \checkmark & \checkmark \\
N-1 Reserve (C6)         & \checkmark & \checkmark & \checkmark
  & \checkmark & \checkmark & \checkmark & \checkmark \\
\texttt{con\_battery}    & ---        & \checkmark & \checkmark
  & \checkmark & \checkmark & \checkmark & \checkmark \\
\texttt{con\_wt}         & ---        & ---        & \checkmark
  & ---        & \checkmark & \checkmark & \checkmark \\
\texttt{con\_min\_inertia} (C12) & ---  & ---      & ---
  & \checkmark & \checkmark & \checkmark & \checkmark \\
\texttt{con\_battery\_vi} & ---       & ---        & ---
  & ---        & ---        & \checkmark & \checkmark \\
\texttt{con\_wt\_vi}     & ---        & ---        & ---
  & ---        & ---        & ---        & \checkmark \\
\texttt{con\_qss\_limit} (C13) & ---  & ---        & ---
  & ---        & ---        & ---        & \checkmark \\
\bottomrule
\multicolumn{8}{l}{\checkmark = Active, --- = Inactive} \\
\end{tabular}
\end{table}

\subsection{Implementation}

The model is implemented in Python~3.9 using the Pyomo~6.9.5
algebraic modeling language and solved with IBM CPLEX~22.1.1
(\texttt{cplex\_direct} interface). Generator, demand, BESS,
and wind turbine data are loaded from
\texttt{DataSet\_ModifikasiCandra\_R02.xlsx}.
All seven scenarios are executed sequentially on a standard
workstation (Intel Core i7, 16~GB RAM, Windows~10).

% =========================================================
\section{Results and Discussion}
\label{sec:results}

\subsection{Overview of All Seven Scenarios}

Table~\ref{tab:results} presents the complete economic and
frequency security performance metrics across all seven scenarios.
Cost values are in USD per 24-hour scheduling day.

\begin{table*}[!t]
\renewcommand{\arraystretch}{1.3}
\caption{Comprehensive Simulation Results Across All Seven Scenarios
(USD/day; 24-Hour Horizon; CPLEX 22.1.1)}
\label{tab:results}
\centering\small
\begin{tabular}{lrrrrrrr}
\toprule
\textbf{Metric} & \textbf{S1} & \textbf{S2} & \textbf{S3}
  & \textbf{S4} & \textbf{S5} & \textbf{S6} & \textbf{S7} \\
\midrule
\multicolumn{8}{l}{\textit{Cost Components (USD/day)}} \\
Fuel Cost         & 3{,}452{,}429 & 3{,}506{,}628 & 3{,}226{,}210
  & 3{,}620{,}477 & 3{,}388{,}277 & 3{,}383{,}101 & 3{,}383{,}209 \\
Fixed Cost        &   514{,}954 &   507{,}545 &   487{,}590
  &   555{,}881 &   555{,}881 &   555{,}881 &   555{,}881 \\
Startup Cost      &    61{,}840 &     5{,}712 &     2{,}592
  &           0 &           0 &           0 &           0 \\
\textbf{Total Cost (USD)} & \textbf{4{,}029{,}223} & \textbf{4{,}019{,}885}
  & \textbf{3{,}716{,}391} & \textbf{4{,}176{,}358} & \textbf{3{,}944{,}158}
  & \textbf{3{,}938{,}983} & \textbf{3{,}939{,}091} \\
\midrule
Cost vs.\ S1 (\%)  & $\pm$0.00 & $-$0.23 & $-$7.76 & $+$3.65
  & $-$2.11 & $-$2.24 & $-$2.24 \\
Cost vs.\ S4 (\%)  & $-$3.52 & $-$3.74 & $-$11.04 & $\pm$0.00
  & $-$5.57 & $-$5.70 & $-$5.69 \\
\midrule
\multicolumn{8}{l}{\textit{System Inertia}} \\
Avg.\ $H_{\text{sys}}$ [MWs] & 36{,}623 & 35{,}844 & 35{,}454
  & 42{,}980 & 42{,}980 & 43{,}216 & 43{,}256 \\
Min.\ $H_{\text{sys}}$ [MWs] & 32{,}565 & 35{,}400 & 35{,}400
  & 42{,}690 & 42{,}690 & 42{,}690 & 42{,}690 \\
Avg.\ BESS-VI [MWs] & 0.0 & 0.0 & 0.0 & 0.0 & 0.0 & 236.5 & 263.5 \\
Avg.\ WT-VI [MWs]   & 0.0 & 0.0 & 0.0 & 0.0 & 0.0 &   0.0 &  12.7 \\
\midrule
\multicolumn{8}{l}{\textit{Frequency Security}} \\
Worst RoCoF [Hz/s]       & 0.903 & 0.831 & 0.831 & 0.550 & 0.550
  & 0.550 & 0.550 \\
RoCoF Viol.\ (of 24~h)   &  24   &  24   &  24   &   0   &   0
  & ${\approx}$0 & ${\approx}$0 \\
\midrule
\multicolumn{8}{l}{\textit{Renewable Integration}} \\
Avg.\ Units ON  & 6.6 & 6.6 & 6.1 & 7.0 & 7.0 & 7.0 & 7.0 \\
WT Curtailment  & --- & --- & 0\% & --- & 1.3\% & 1.3\% & 1.3\% \\
\bottomrule
\end{tabular}
\end{table*}

\subsection{Baseline Classical UC Without Frequency Constraints (S1)}

Scenario~1 establishes the cost baseline at \$4{,}029{,}223/day.
The worst-case estimated RoCoF is \textbf{0.903~Hz/s}---64\% above
the 0.55~Hz/s limit---with violations in \textbf{all 24} scheduling
periods. Average system inertia is only 36{,}623~MWs, with a
minimum of 32{,}565~MWs during off-peak hours. At this inertia level,
the loss of the largest committed unit causes RoCoF that would
trigger relay operation in practical systems. This confirms that
classical UC is frequency-insecure under N-1 contingency conditions.

\subsection{BESS and Wind Turbine Integration Without Frequency Constraints (S2, S3)}

\textbf{Scenario~2} (classical UC + BESS) achieves a marginal
cost reduction of only 0.23\% versus S1 (\$4{,}019{,}885/day).
The BESS reduces startup costs from \$61{,}840 to \$5{,}712 by
smoothing commitment scheduling, but slightly increases fuel cost.
Frequency security remains violated in all 24 hours.

\textbf{Scenario~3} (classical UC + BESS + WT) achieves the
\emph{greatest single cost reduction} of \textbf{7.76\%}
versus S1, reaching \$3{,}716{,}391/day. Wind generation
(400~MW) reduces thermal fuel consumption by \$226{,}220/day.
Startup costs fall to \$2{,}592. Zero wind curtailment is
observed, indicating full wind utilization. However,
frequency security remains violated in all 24 hours.

\subsection{Enforcing RoCoF Constraint via SG Inertia Only (S4, S5)}

\textbf{Scenario~4} (FCUC + BESS + minimum inertia constraint only,
no WT) requires additional must-run thermal generators to satisfy
constraint~\eqref{eq:rocof_con}, increasing total cost to
\$4{,}176{,}358/day ($+$3.65\% vs.\ S1).
Fixed costs rise from \$514{,}954 to \$555{,}881 due to
additional committed units, while startup costs vanish (units
run continuously).
All 24 hours satisfy RoCoF~$\leq$~0.55~Hz/s with an
average $H_{\text{sys}}$ of 42{,}980~MWs.

\textbf{Scenario~5} (FCUC + BESS + WT + minimum inertia) demonstrates
that adding wind generation \emph{while enforcing} frequency security
recovers significant cost: \$3{,}944{,}158/day ($-$5.57\% vs.\ S4,
$-$2.11\% vs.\ S1). Wind reduces fuel cost by \$232{,}200/day
versus S4. Minor wind curtailment of 1.3\% appears as the optimizer
balances wind output against reserve and inertia requirements.

\subsection{Virtual Inertia Contribution (S6, S7)}

\textbf{Scenario~6} activates BESS virtual inertia
(\texttt{con\_battery\_vi}), contributing an average of
\textbf{236.5~MWs} additional inertia via grid-forming converter
control. Total cost: \$3{,}938{,}983/day ($-$5.70\% vs.\ S4).

\textbf{Scenario~7} (the proposed full model) further activates WT
virtual inertia (\texttt{con\_wt\_vi}) and QSS constraint
(\texttt{con\_qss\_limit}). BESS-VI averages \textbf{263.5~MWs}
and WT-VI averages \textbf{12.7~MWs} of additional inertia.
Total $H_{\text{sys,avg}}$ = 43{,}256~MWs.
Total cost: \$3{,}939{,}091/day---essentially identical to
S6~(+\$108, ${<}0.003\%$).

The marginal cost of enabling WT-VI provision is negligible
(\$108/day), while it provides additional frequency security margin.
WT-VI operates within the 1.3\% curtailment headroom without
increasing curtailment.

\subsection{Economic Analysis: Cost-Security Trade-off}

Table~\ref{tab:econ} summarizes the economic efficiency of
frequency security enforcement.

\begin{table}[!t]
\renewcommand{\arraystretch}{1.2}
\caption{Economic Efficiency of Frequency Security Enforcement}
\label{tab:econ}
\centering\footnotesize
\begin{tabular}{lrcc}
\toprule
Scenario & Total Cost & vs.\ S1 & Freq.\ Secure? \\
         & [USD/day]  & [\%]    & \\
\midrule
S1: Classical UC           & 4{,}029{,}223 & $\pm$0.00 & No (24/24) \\
S4: FCUC, SG-inertia only  & 4{,}176{,}358 & $+$3.65   & Yes \\
S5: FCUC + WT              & 3{,}944{,}158 & $-$2.11   & Yes \\
S6: FCUC + BESS-VI         & 3{,}938{,}983 & $-$2.24   & Yes \\
\textbf{S7: Proposed (Full VI)} & \textbf{3{,}939{,}091} & $\mathbf{-}$\textbf{2.24}
  & \textbf{Yes} \\
\bottomrule
\end{tabular}
\end{table}

\textbf{Key findings:}
\begin{enumerate}
  \item The proposed multi-source FCUC (S7) achieves frequency
    security at \emph{lower total cost than the classical UC
    baseline} ($-$2.24\%), demonstrating that co-optimizing
    renewable integration and virtual inertia does not require
    cost penalties.

  \item Multi-source VI (S7 vs.\ S4) saves \textbf{\$237{,}267/day}
    (5.69\%), primarily through reduced must-run thermal commitment.

  \item Adding WT-VI to BESS-VI (S7 vs.\ S6) costs only \$108/day
    while providing additional inertia (+40~MWs average), indicating
    WT-VI is effectively ``free'' ancillary service.
\end{enumerate}

% =========================================================
\section{Conclusion}
\label{sec:conclusion}

This paper proposed, implemented, and validated a MILP-based
Frequency-Constrained Unit Commitment model incorporating
multi-source virtual inertia from BESS and wind turbines,
solved using Python~3.9/Pyomo and IBM CPLEX~22.1.1.
Seven simulation scenarios on a modified IEEE 10-unit test system
yielded the following conclusions:

\begin{enumerate}
  \item \textbf{Classical UC is frequency-insecure.} Scenario~1
    shows a worst-case RoCoF of 0.903~Hz/s (limit: 0.55~Hz/s)
    violated in all 24 hours, confirming the necessity of FCUC.

  \item \textbf{Wind integration provides the greatest cost reduction.}
    Scenario~3 achieves $-$7.76\% versus S1 through fuel cost savings
    of \$226{,}220/day, but does not address frequency security.

  \item \textbf{RoCoF enforcement via SG inertia alone is costly.}
    Scenario~4 incurs $+$3.65\% versus S1, requiring more must-run
    thermal units.

  \item \textbf{Multi-source VI recovers the frequency security premium.}
    Scenario~7 saves 5.69\% versus S4 (\$237{,}267/day) while
    maintaining full RoCoF compliance; average system inertia rises
    to 43{,}256~MWs (263.5~MWs from BESS-VI, 12.7~MWs from WT-VI).

  \item \textbf{Frequency security achievable below baseline cost.}
    Scenario~7 costs 2.24\% \emph{less} than classical UC while
    being fully frequency-secure, demonstrating that co-optimizing
    renewable integration and VI provision simultaneously reduces
    costs and improves grid security.
\end{enumerate}

For Indonesia's RUPTL 2025--2034 goal of 61\% renewable generation,
these results indicate that explicit scheduling of virtual inertia
from planned BESS and wind assets can fulfill inertia requirements
without additional cost burden.
Future work will extend the model to: (i) stochastic formulation
for renewable forecast uncertainty; (ii) transmission-constrained
(DC power flow) FCUC; and (iii) validation on actual Indonesian
power system topology (JAMALI or Sulawesi grid).

% =========================================================
\bibliographystyle{IEEEtran}
\begin{thebibliography}{20}

\bibitem{PLN2025}
Perusahaan Listrik Negara (PLN),
``Rencana Usaha Penyediaan Tenaga Listrik (RUPTL) PLN
2025--2034,'' PT PLN (Persero), Jakarta, Indonesia, 2025.
[Online]. Available: \url{https://www.pln.co.id/ruptl}

\bibitem{Tielens2016}
P.~Tielens and D.~Van~Hertem,
``The relevance of inertia in power systems,''
\textit{Renew. Sustain. Energy Rev.},
vol.~55, pp.~999--1009, Mar.~2016.
DOI: \href{https://doi.org/10.1016/j.rser.2015.11.016}
{10.1016/j.rser.2015.11.016}

\bibitem{AEMO2017}
Australian Energy Market Operator (AEMO),
``Black System South Australia 28 September 2016---Final Report,''
AEMO, Melbourne, Australia, 2017.
[Online]. Available:
\url{https://www.aemo.com.au/Electricity/National-Electricity-Market-NEM/
Market-notices-and-events/Power-System-Incident-Reports/2017/
Integrated-Final-Report-SA-Black-System-28-September-2016}

\bibitem{E02Sulawesi}
L.~M.~Putranto \textit{et al.},
``Unit Commitment with Primary Frequency Regulation Consideration
in the Southern Sulawesi Power System,''
in \textit{Proc. 2024 Int. Conf. on ICT for Power and Energy
(ICT-PEP)}, 2024.

\bibitem{E03JAMALI}
L.~M.~Putranto \textit{et al.},
``Unit Commitment with 2-Minute Fast Frequency Reserve Constraint
in the Jawa-Madura-Bali Power System,''
in \textit{Proc. 2025 Int. Conf. on ICT for Power and Energy
(ICT-PEP)}, 2025.

\bibitem{Padhy2004}
N.~P.~Padhy,
``Unit commitment---a bibliographical survey,''
\textit{IEEE Trans. Power Syst.},
vol.~19, no.~2, pp.~1196--1205, May~2004.
DOI: \href{https://doi.org/10.1109/TPWRS.2003.821611}
{10.1109/TPWRS.2003.821611}

\bibitem{Conejo2010}
A.~J.~Conejo, M.~Carri\'on, and J.~M.~Morales,
\textit{Decision Making Under Uncertainty in Electricity Markets}.
New York, NY, USA: Springer, 2010.
DOI: \href{https://doi.org/10.1007/978-1-4419-7421-1}
{10.1007/978-1-4419-7421-1}

\bibitem{Wen2016}
Y.~Wen, W.~Li, G.~Huang, and X.~Liu,
``Frequency dynamics constrained unit commitment with battery
energy storage,''
\textit{IEEE Trans. Power Syst.},
vol.~31, no.~6, pp.~5115--5125, Nov.~2016.
DOI: \href{https://doi.org/10.1109/TPWRS.2016.2521882}
{10.1109/TPWRS.2016.2521882}

\bibitem{Anderson1990}
P.~M.~Anderson and M.~Mirheydar,
``A low-order system frequency response model,''
\textit{IEEE Trans. Power Syst.},
vol.~5, no.~3, pp.~720--729, Aug.~1990.
DOI: \href{https://doi.org/10.1109/59.65898}
{10.1109/59.65898}

\bibitem{Ahmadi2014}
H.~Ahmadi and H.~Ghasemi,
``Security-constrained unit commitment with linearized system
frequency limit constraints,''
\textit{IEEE Trans. Power Syst.},
vol.~29, no.~4, pp.~1536--1545, Jul.~2014.
DOI: \href{https://doi.org/10.1109/TPWRS.2014.2297997}
{10.1109/TPWRS.2014.2297997}

\bibitem{Teng2016}
F.~Teng, V.~Trovato, and G.~Strbac,
``Stochastic scheduling with inertia-dependent fast frequency
response requirements,''
\textit{IEEE Trans. Power Syst.},
vol.~31, no.~2, pp.~1557--1566, Mar.~2016.
DOI: \href{https://doi.org/10.1109/TPWRS.2015.2434837}
{10.1109/TPWRS.2015.2434837}

\bibitem{Chu2020}
Z.~Chu, U.~Markovic, G.~Hug, and F.~Teng,
``Towards optimal system scheduling with synthetic inertia
provision from wind turbines,''
\textit{IEEE Trans. Power Syst.},
vol.~35, no.~5, pp.~4056--4066, Sep.~2020.
DOI: \href{https://doi.org/10.1109/TPWRS.2020.2985843}
{10.1109/TPWRS.2020.2985843}

\bibitem{Shi2025}
Q.~Shi, Y.~Zhu, K.~Fan, R.~Cheng, J.~Shen, and W.~Liu,
``Two-stage linearization of frequency nadir constraint for
unit commitment,''
\textit{IEEE Trans. Power Syst.},
vol.~40, no.~4, pp.~3584--3587, Jul.~2025.
DOI: \href{https://doi.org/10.1109/TPWRS.2025.3564110}
{10.1109/TPWRS.2025.3564110}

\bibitem{Liu2024}
X.~Liu, X.~Fang, N.~Gao, H.~Yuan, A.~Hoke, H.~Wu, and J.~Tan,
``Frequency nadir constrained unit commitment for high renewable
penetration island power systems,''
\textit{IEEE Open Access J. Power Energy},
vol.~11, pp.~141--153, 2024.
DOI: \href{https://doi.org/10.1109/OAJPE.2024.3370504}
{10.1109/OAJPE.2024.3370504}

\bibitem{Asiedu2025}
S.~T.~Asiedu, T.~Aryal, A.~Suvedi, Z.~Wang,
H.~M.~Rekabdarkolaee, and T.~M.~Hansen,
``A review of dynamic constraint integration into steady-state
power system optimization models for IBR-dominated systems,''
\textit{IEEE Access},
vol.~13, pp.~202665--202699, 2025.
DOI: \href{https://doi.org/10.1109/ACCESS.2025.3636078}
{10.1109/ACCESS.2025.3636078}

\bibitem{VI06OptimalPlacement}
H.~Fang, L.~Chen, N.~Dlakic, and D.~Ling,
``Optimal energy storage system-based virtual inertia placement:
A frequency stability point of view,''
\textit{IEEE Trans. Power Syst.}, vol.~36, no.~3,
pp.~2324--2337, May~2021.
DOI: \href{https://doi.org/10.1109/TPWRS.2021.3052789}
{10.1109/TPWRS.2021.3052789}

\bibitem{Kundur1994}
P.~Kundur,
\textit{Power System Stability and Control}.
New York, NY, USA: McGraw-Hill, 1994.
ISBN: 978-0-07-035958-1

\bibitem{Kazarlis1996}
S.~A.~Kazarlis, A.~G.~Bakirtzis, and V.~Petridis,
``A genetic algorithm solution to the unit commitment problem,''
\textit{IEEE Trans. Power Syst.},
vol.~11, no.~1, pp.~83--92, Feb.~1996.
DOI: \href{https://doi.org/10.1109/59.485989}
{10.1109/59.485989}

\bibitem{Carrion2006}
M.~Carri\'on and J.~M.~Arroyo,
``A computationally efficient mixed-integer linear formulation
for the thermal unit commitment problem,''
\textit{IEEE Trans. Power Syst.},
vol.~21, no.~3, pp.~1371--1378, Aug.~2006.
DOI: \href{https://doi.org/10.1109/TPWRS.2006.876672}
{10.1109/TPWRS.2006.876672}

\bibitem{Xu2018}
T.~Xu, W.~Jang, and T.~J.~Overbye,
``Commitment of fast-responding storage devices to mimic inertia
for the enhancement of primary frequency response,''
\textit{IEEE Trans. Power Syst.},
vol.~33, no.~2, pp.~1219--1230, Mar.~2018.
DOI: \href{https://doi.org/10.1109/TPWRS.2017.2735990}
{10.1109/TPWRS.2017.2735990}

\bibitem{D02ForecastError}
N.~Nguyen-Hong and Y.~Nakanishi,
``Frequency-constrained unit commitment considering battery
storage system and forecast error,''
in \textit{Proc. 2018 IEEE Innov. Smart Grid Technol.---Asia
(ISGT Asia)}, Singapore, 2018, pp.~290--295.
DOI: \href{https://doi.org/10.1109/ISGT-Asia.2018.8467848}
{10.1109/ISGT-Asia.2018.8467848}

\end{thebibliography}

\end{document}
